\documentclass{ibetest}
\renewcommand{\arraystretch}{1.35}
\begin{document}
\testheader
\vspace{16mm}
\begin{center}
  {\sffamily\bfseries\Huge Thermo-S1 \,$\cdot$\, Part 1}\\[6pt]
  {\sffamily\bfseries\LARGE Progress Tests 1.A $\to$ 3.C}\\[8pt]
  {\large Description of a system in thermodynamics}\\[3pt]
  {\normalsize One short test per lecture section, in the format of Test~1}\\[3pt]
  {\normalsize Semester 1 \,$\cdot$\, Academic Year 2026--2027}
\end{center}
\vspace{10mm}

{\sffamily\bfseries How to use this booklet}
\begin{enumerate}[leftmargin=7mm, itemsep=3pt]
  \item Study a section of the lecture notes, then take its progress test \emph{in real
        conditions}: timer on, no notes, no calculator, answers only inside the frames.
  \item Mark yourself with the answer-key booklet. Each circled number in the margin is one
        marking item; the scheme below each exercise says exactly what earns it.
  \item Convert your score: mark out of 20 $=$ score $\times\,20/14$. Write it in the tracker
        (next page). Below 14/20, re-read the section and retake the test a few days later.
  \item When all nine tests are done, take the \textbf{End-of-Unit Test} (45 minutes,
        50 marks, sections 1.A to 3.C).
\end{enumerate}

\vspace{6mm}
{\sffamily\bfseries Contents}\par\vspace{2mm}
\begin{center}
\begin{tabular}{@{}l l c c c@{}}
\hline
\textbf{Test} & \textbf{Lecture section} & \textbf{Notes, pages} & \textbf{Duration} &
\textbf{Marks}\\ \hline
1.A & System and state variables                      & 2--5   & 7 min  & 14\\
1.B & Postulates of thermodynamics                    & 5--7   & 7 min  & 14\\
1.C & Equations of state and state functions          & 7--10  & 7 min  & 14\\
2.A & Thermoelastic behaviour of condensed phases     & 11--13 & 10 min & 14\\
2.B & Thermoelastic behaviour of gases                & 13--14 & 7 min  & 14\\
2.C & Thermoelastic coefficients: principle and method & 14--15 & 10 min & 14\\
3.A & Low-density gases: the historical laws          & 15--17 & 7 min  & 14\\
3.B & The ideal gas and Clapeyron's law               & 17--20 & 10 min & 14\\
3.C & Equations of state of real gases                & 20--23 & 10 min & 14\\ \hline
\end{tabular}
\end{center}

\vspace{6mm}
\begin{tcolorbox}[enhanced, colback=black!3, colframe=black!45, boxrule=0.4pt, arc=1mm,
  fontupper=\small]
Every test mirrors Test~1 (21 September 2026): a unit conversion, a definition with its
equation and SI units, tick-box questions, and a short calculation with the data framed. The
data use $R = 8.3\ \mathrm{J\,K^{-1}\,mol^{-1}}$ and $RT\approx2.5\times10^{3}\
\mathrm{J\,mol^{-1}}$ at 300~K, so that everything can be done without a calculator.
\end{tcolorbox}

\newpage
\testheader
\vspace{8mm}
{\sffamily\bfseries\Large Progress tracker}\par\vspace{4mm}
\begin{center}
\renewcommand{\arraystretch}{2.1}
\begin{tabular}{|c|p{22mm}|c|c|p{22mm}|c|p{42mm}|}
\hline
\textbf{Test} & \textbf{Date} & \textbf{/14} & \textbf{/20} & \textbf{Retake date} &
\textbf{/14} & \textbf{To revise}\\ \hline
1.A & & & & & & \\ \hline
1.B & & & & & & \\ \hline
1.C & & & & & & \\ \hline
2.A & & & & & & \\ \hline
2.B & & & & & & \\ \hline
2.C & & & & & & \\ \hline
3.A & & & & & & \\ \hline
3.B & & & & & & \\ \hline
3.C & & & & & & \\ \hline\hline
\multicolumn{2}{|l|}{\textbf{End of unit}} & \multicolumn{2}{c|}{\hspace{6mm}/50} &
\multicolumn{3}{l|}{\quad /20 $=$ score $\times\,20/50$} \\ \hline
\end{tabular}
\end{center}
\vspace{6mm}
{\sffamily\bfseries Before every test, remember Test~1}
\begin{itemize}[leftmargin=7mm, itemsep=2pt]
  \item Absolute temperatures in kelvins, $T = \theta + 273.15$, written with a capital K
        and no degree sign. A temperature \emph{difference} is the same in K and in $^\circ$C.
  \item SI units before using $pV = nRT$: $1\ \mathrm L = 10^{-3}\ \mathrm{m^3}$,
        $1\ \mathrm{cm^2} = 10^{-4}\ \mathrm{m^2}$, $1\ \mathrm{bar} = 10^{5}\ \mathrm{Pa}$,
        molar masses in $\mathrm{kg\,mol^{-1}}$.
  \item Letters first, numbers second, with units; powers of ten separately from the mantissas.
  \item Before choosing a law, list which variables are fixed and which are free.
\end{itemize}
\end{document}
